\newpage
\section*{List of Publications} % Use \section* instead of \chapter*
\addcontentsline{toc}{chapter}{List of Publications} % Add to ToC at chapter level
\markboth{Summary}{} % Set header text

\vspace{0.3cm} % Adjust space after title

This doctoral thesis has been written following the Compendium of Articles format and comprises 3 published articles and 2 unpublished manuscripts divided into three chapters. Below, the list of publications included in this thesis is presented. For each article, full bibliographic information, co-authorship details, journal quality metrics, and a description of the author's contributions are provided. Journal quality metrics correspond to the impact factor and quartile for the year of publication, as reported in the Journal Citation Reports (JCR), as provided by the Fundación Española para la Ciencia y la Tecnología (FECYT).

Chapter \ref{chap1} \textbf{Targeting the CAIn Landscape: From Inflammation to Membrane-Specific Metadynamics} contains two articles that establish the biological framework of the Cancer-Aging-Infection (CAIn) triangle and quantify the specificity of antimicrobial peptide interactions with pathological membranes using enhanced sampling methods:
\begin{enumerate}
    \item \textbf{Authors:} \textbf{Daniel Conde-Torres}, Alexandre Blanco-González, Alejandro Seco-González, Fabián Suárez-Lestón, Alfonso Cabezón, Paula Antelo-Riveiro, Ángel Piñeiro, and Rebeca García-Fandiño.
    
    \textbf{Citation:} Conde-Torres, D., et al. Unraveling lipid and inflammation interplay in cancer, aging and infection for novel theranostic approaches. Frontiers in Immunology 15, 1320779 (2024). doi: \href{https://doi.org/10.3389/fimmu.2024.1320779}{10.3389/fimmu.2024.1320779}.

    \textbf{Metrics:} In 2024, Frontiers in Immunology was ranked Q1 in JCR, with an Impact Factor of 5.9.

    \textbf{Contributions:} As first author of the article, I contributed to data curation, investigation, writing the original draft, and reviewing and editing the manuscript.

    \item \textbf{Authors:} \textbf{Daniel Conde-Torres}, Martín Calvelo, Carme Rovira, Ángel Piñeiro, and Rebeca Garcia-Fandino.
    
    \textbf{Citation:} Conde-Torres, D., Calvelo, M., Rovira, C., Piñeiro, Á. \& Garcia-Fandino, R. Unlocking the specificity of antimicrobial peptide interactions for membrane-targeted therapies. Computational and Structural Biotechnology Journal 25, 61–74 (2024). doi: \href{https://doi.org/10.1016/j.csbj.2024.04.022}{10.1016/j.csbj.2024.04.022}.

    \textbf{Metrics:} In 2024, the Computational and Structural Biotechnology Journal was ranked Q2 in JCR, with an Impact Factor of 4.1.

    \textbf{Contributions:} As first author of the article, I contributed to visualization, validation, methodology, investigation, and writing the original draft.
\end{enumerate}

Chapter \ref{chap2} \textbf{Transcending Classical Limits: Quantum Computing for Peptide Folding} is formed by one article that introduces a resource-efficient quantum algorithm to model peptide folding at membrane boundaries:
\begin{enumerate}
    \item \textbf{Authors:} \textbf{Daniel Conde-Torres}, Mariamo Mussa-Juane, Daniel Faílde, Andrés Gómez, Rebeca García-Fandiño, and Ángel Piñeiro.

    \textbf{Citation:} Conde-Torres, D., et al. Classical Simulations on Quantum Computers: Interface-Driven Peptide Folding on Simulated Membrane Surfaces. Computers in Biology and Medicine 182, 109157 (2024). doi: \href{https://doi.org/10.1016/j.compbiomed.2024.109157}{10.1016/j.compbiomed.2024.109157}.

    \textbf{Metrics:} In 2024, Computers in Biology and Medicine was ranked Q1 in JCR, with an Impact Factor of 6.3.

    \textbf{Contributions:} As first author of the article, I contributed to the software development, methodology, formal analysis, data curation, validation, visualization, and writing the original draft.
\end{enumerate}

Chapter \ref{chap3} \textbf{Environment-Conditioned Design: A Physics-Based Hamiltonian Pivot} is formed by two unpublished manuscripts detailing a novel approach for the automated sequence design of membrane-interacting peptides and its deployment as an open-source web platform:
\begin{enumerate}
    \item \textbf{Authors:} \textbf{Daniel Conde-Torres}, Rebeca García-Fandiño, and Ángel Piñeiro.

    \textbf{Citation:} Conde-Torres, D., García-Fandiño, R. \& Piñeiro, Á. Environment-conditioned design of $\alpha$-helical peptides. bioRxiv preprint (2026). doi: \href{https://doi.org/10.1101/2026.05.07.723485}{10.1101/2026.05.07.723485}.

    \textbf{Metrics:} None (unpublished manuscript).

    \textbf{Contributions:} As first author of the article, I contributed to the conceptualization of the project, developed the Hamiltonian framework, conducted the computational optimization experiments across the quantum and classical backends, and contributed to writing the original draft.

    \item \textbf{Authors:} \textbf{Daniel Conde-Torres}, César Porto, Rebeca García-Fandiño, and Ángel Piñeiro.

    \textbf{Citation:} Conde-Torres, D., Porto, C., García-Fandiño, R. \& Piñeiro, Á. Helix Design Studio: A Web Platform for Environment-Conditioned Design of $\alpha$-Helical Peptides. bioRxiv preprint (2026).

    \textbf{Metrics:} None (unpublished manuscript).

    \textbf{Contributions:} As first author of the article, I contributed to the conceptualization of the platform, participated in the full-stack software development (including the integration of the Hamiltonian algorithm and the 3D molecular viewer), and contributed to writing the results[cite: 5].
\end{enumerate}